\documentclass[11pt]{article}
\usepackage[a4paper,margin=25mm]{geometry}
\usepackage{amsmath,amssymb,amsthm,booktabs,longtable,hyperref,xurl,microtype}
\setlength{\emergencystretch}{3em}
\hypersetup{hidelinks}
\newtheorem{theorem}{Theorem}
\newtheorem{lemma}[theorem]{Lemma}
\newtheorem{proposition}[theorem]{Proposition}
\newtheorem{corollary}[theorem]{Corollary}
\newcommand{\Q}{\mathbb Q}\newcommand{\Z}{\mathbb Z}
\title{One contact gauge in a combined quadratic system\\
Quartic checkpoints in 1185 and 1355 variables}
\author{\(\mathbb Z\mathbb Q\) Diophantine Project}
\date{23 November 2026}
\begin{document}
\maketitle
\begin{abstract}
We construct complete rational quadratic systems with 1185 variables
and 1095 equations, or 1355 variables and 1241 equations. Their sums of
squares have degree exactly four. The starting systems have 1601 and
1837 variables and contain twelve integrality certificates. We replace
their twelve contact gauges by the sum of all contact weights. This
preserves the original blockwise valuation argument and its height
thresholds, while allowing 91 or 93 stored coordinates to be shared.
Reordering the three common detector factors saves their intermediate
products, and an exact finite-pair identity removes one private field
from each input chart at the cost of two common fields.

We supply the full arithmetic transfer under the retained interfaces,
the physical-function dictionary, exact reconstruction maps, and a
complete independent coefficient verifier. The universality assertion
\(U_{\mathrm{em}}\) for the frozen emended integer source remains open.
The work therefore supplies conditional quartic reductions and exact
algebraic checkpoints, not an unconditional undecidability theorem.
\end{abstract}

\section{Scope, hypotheses, and the fixed starting system}

The predecessor systems \cite{common} have 1601 variables and 1599
quadratic equations in the general interface, or 1837 variables and
1811 equations in the specified interface. Both use the same optimized
outer circuit, with 251 rational coordinates and 219 quadratic
equations. Eleven ordinary integer inputs and one reciprocal integer
input are certified by the 163-coordinate or 179-coordinate
integer-defining interface. Those counts include the distinguished
input of a defining block.

The outer circuit in this paper is frozen exactly as in the predecessor
tables. In particular, this paper does not repair or assert universality
of its underlying emended positive-integer source. We denote that
unresolved assertion by \(U_{\mathrm{em}}\), as in the source
reconciliation of \cite{outer}. The consulted source displays have
incompatible coefficients and Pell-row typography; the exact published
universality theorem for the frozen emendation has not been verified.
Any later correction to that integer source must receive its own
reconstruction proof. The coefficient tables here record
\texttt{source\_universality\_verified: false} and
\texttt{unconditional\_undecidability\_claim: false}.

The retained arithmetic interfaces are those stated in
\cite{gate,scalar,common}: the group and formal-logarithm descriptions on
the specified elliptic curve over \(K=\Q(\sqrt2)\); ordered scalar
extraction and preparation; full-domain reflected-contact and positive
aggregate height estimates; the native-pole obstruction and prepared
realization theorem; and the finite-family Gaussian prime supplier,
including its stated recurrence and all-sign order coprimality. The
simultaneous common-point construction of \cite{common,crt} is used
without changing the private input differences. The general route
retains its sufficiently large fixed defining parameter, whose
threshold is ineffective. The specified route retains its already
specified parameter and estimates. No ordinary rational \(abc\)
assumption is added by the present transfer.

All comparisons of rational zero sets below are under these explicit
interfaces. A finite coefficient check verifies polynomial identities,
counts, and maps; it does not establish the retained arithmetic
premises or \(U_{\mathrm{em}}\).

\section{One contact gauge for all input blocks}
\label{sec:global-contact-gauge}

Let the input blocks be indexed by \(1\le i\le k\), and let
\(0\le j<e\) index their contact arguments, where \(e=2\)
for the general interface and \(e=3\) for the specified
interface. Retain every block's original contact function
\(q_{ij}\), original contact argument, and original aggregate
relation. We make the single change
\begin{equation}\label{eq:global-contact-sum}
                    c=\sum_{i=1}^{k}\sum_{j=0}^{e-1}q_{ij},
                    \qquad c\ne0,
\end{equation}
in place of the \(k\) separate contact sums. Every product
\(W_i\in K=\mathbb Q(\sqrt2)\), every rational norm/trace
error \(D_i\), and every native-pole witness remains private
to its input block. Define
\begin{equation}\label{eq:global-gates}
 \begin{split}
 D_i\operatorname{Tr}(W_i)&=\operatorname{N}(W_i-1),\\
 s_{ij}&=c-q_{ij},\qquad
 g_{ij}=1+D_i^2+\frac{s_{ij}^2}{q_{ij}} .
 \end{split}
\end{equation}
The actual comparison-pole arguments remain
\begin{equation}\label{eq:global-pole}
               b_{ij}=(408/409)^2 R_{ij}g_{ij}^2/4,
\end{equation}
with precisely the source contact functions \(R_{ij}\)
and source cleared equations. No new division or
nonvanishing condition at an old boundary is introduced.

The original proof of the shared gate uses the equality
\(c=q_{ij}+s_{ij}\), rather than the number of summands in
the definition of \(c\). We give the entire valuation
argument because projecting the new system to an old
block does not restore that block's former contact-sum
identity.

\begin{lemma}[The gate with an arbitrary common gauge]
\label{lem:global-gate}
Fix an input block \(i\), one of its consumed primes
\(p\), and a contact \(j\) with
\[
            v_p(q_{ij})=-\mu,\qquad
            v_p(R_{ij})=\mu,\qquad \mu>0\ \hbox{odd}.
\]
Retain the source pole obstruction and the source
exceptional-prime exclusions. Every rational solution
of \eqref{eq:global-gates}--\eqref{eq:global-pole}
satisfies
\begin{equation}\label{eq:global-gate-bounds}
 v_p(D_i)>-\mu/4,\qquad
 v_p(s_{ij})>-3\mu/4,\qquad v_p(c)=-\mu.
\end{equation}
If every factor of \(W_i\) has nonpositive valuation
at both conjugate slots above \(p\), then one slot has
\[
                              v(W_i)>-\mu/4 .
\]
No property of the other input blocks is required
for this conclusion.
\end{lemma}
\begin{proof}
First, the norm/trace equation forces
\(\operatorname{Tr}(W_i)\ne0\). Otherwise
\(\operatorname{N}(W_i-1)=0\), and the field property
would give \(W_i=1\), whose trace is 2. The same
equation rules out \(W_i=0\), since its right side
would be 1. Thus \(D_i\) is the actual rational
norm/trace quotient.

Put \(v(0)=+\infty\). If \(g_{ij}\ne0\), then
\eqref{eq:global-pole} has odd valuation
\[
       \mu+2v_p(g_{ij})+2v_p(408/409).
\]
Outside the source exceptional set,
\(v_p(408/409)\le0\). This includes \(p=409\),
where its negative value strengthens the
obstruction. Since a negative odd pole valuation
is forbidden, \(v_p(g_{ij})>-\mu/2\).
The same inequality is automatic if \(g_{ij}=0\).

The nonconstant summands of \(g_{ij}\) have
valuations \(2v_p(D_i)\), which is even, and
\(2v_p(s_{ij})+\mu\), which is odd.
A negative minimum is therefore unique and
cannot cancel with the other summand or the
constant term. Both summands must have valuation
greater than \(-\mu/2\). This proves the first
two inequalities in \eqref{eq:global-gate-bounds}.
Since \(v_p(s_{ij})>-\mu\), the equality
\(c=q_{ij}+s_{ij}\) gives \(v_p(c)=-\mu\).
This argument does not use
\eqref{eq:global-contact-sum}.

At the two conjugate slots, counted with
repetition in the inert case, let \(w_1,w_2\)
be the valuations of \(W_i\). They are
nonpositive by the hypothesis. If either is
zero the conclusion holds. Otherwise both
are negative, so
\[
 v_p\operatorname{N}(W_i-1)=w_1+w_2,\qquad
 v_p\operatorname{Tr}(W_i)\ge\min(w_1,w_2).
\]
Consequently
\[
           v_p(D_i)\le\max(w_1,w_2).
\]
The already proved lower bound on \(v_p(D_i)\)
gives the required slot.
\end{proof}

\subsection{Soundness and the unchanged height thresholds}

Retain all four field-factor polynomials of the v82
interface, with their enhanced powers of \(c\).
For \(v_p(c)=-\mu\), \(\mu\) odd, the source
residue-class argument gives a unique minimum in
each factor, a nonpositive factor valuation, and
residue one for a factor of valuation zero.
These statements depend on the value of \(c\)
and on the unchanged point relations; they
do not depend on how many contact functions
were added to form \(c\).

Lemma~\ref{lem:global-gate} therefore supplies,
for every consumed contact of every input block,
the same successful slot as in v82. Every
factor and every one of its terms has valuation
greater than \(-\mu/4\). The original reciprocal,
cleared-comparison, input, and weighted scalar
estimates follow with exactly their former
constants. In particular, each block recovers
its original integer-label divisibilities
\[
       M_{ij}^{\,3}\mid a\Delta,\qquad
       M_{ij}^{\,3}\mid
       \left|\operatorname{num}(\beta_i-n_i)\right|
\]
when the corresponding residual is nonzero,
with the same treatment of radical one,
zero alternate ordinates, and finite two-torsion
boundaries as in the source.

The later height argument is unchanged
block by block. Its quantities are the block's
scalar normalizer \(R_i\), input \(\beta_i\),
four actual point differences, contact
arguments and radicals. In the general route
the decisive bound is
\[
 B_i^2\le A\bigl(h(R_i)+h(T_i)\bigr)+J,
 \qquad
 B_i=\max(1,|n_i|,|a|,|t|,|b|),
\]
followed by the source contact-height estimates.
There is no bound on \(h(c)\) in that argument.
The specified route likewise uses its original
aggregate shifts, anchor, and three-contact
height constants. Those expressions have not
been changed. Thus neither the sufficient
general-family threshold nor the certified
specified endpoint is enlarged by \(k\).
Every distinguished rational input is integral.

This soundness proof makes no disjointness
assumption on contact radicals of an arbitrary
rational solution. The disjointness used next
belongs only to completeness.

\paragraph{The midpoint boundary \(q_{ij}=-1\).}
All original cleared contact and pole rows are
retained. If \(q_{ij}=-1\), it has no negative
odd valuation at any prime, so its consumed
radical is one. In the reflected source
contact equations this is the midpoint fiber,
which is explicitly included in the retained
full-domain contact-height estimate. The
associated radical divisibilities are vacuous.
The pole quotient involving \(1+q_{ij}\) is
used only at consumed primes, where
\(v_p(q_{ij})<0\) forces \(1+q_{ij}\ne0\).
At the midpoint itself we use the cleared
rows, not that quotient. Those rows may
force the cleared gate value to zero and
leave the pole argument free; no step of
the soundness argument above requires a
quotient on this fiber.

\subsection{Ordered simultaneous completeness}

Fix an arbitrary tuple of integer inputs.
Use the previously proved simultaneous scalar
preparation and common-index construction:
choose the private input root points, common
\(H\), common scalar and Lucas point roles,
and the common input numerator. Keep the
scalar normalizers and scalar-pole witnesses
private. Choose all the finite and real
neighborhoods on which the relevant physical
tests are bounded. The point differences
and their aggregate shifts are fixed at
this stage.

Translate the finite contact lists as in
the common-point proof, so that lists from
different input blocks are disjoint.
The different lists \emph{within} one input
block are also disjoint here. In the general
quartic source their fixed translation
difference is \(T_i>7\), while the reflected
contact set has diameter four. In the
specified three-argument source the two
successive shifts have the retained lower
bounds seven and nine, while the contact
diameter is \(25/11\). These are the
integer-witness separation bounds in
\nolinkurl{papers/quartic/shared_product_audit_20261119_v78.tex},
Section \emph{Integer completeness and the
order of choices}; the v82 scalar-point
proof retains exactly those shifts.
Thus no selected contact in this quartic
construction is a pole of a second weight,
whether from the same input or another.
The same finite-family prime supplier
applies to their union. Enlarge fixed
avoidance to include the values at every
contact of the denominators and other
contact functions, all fixed nonzero
root multipliers, and all earlier
preparation moduli. The common-index
CRT and high-precision approximation
preserve the source synchronized formal
depths and all original local inequalities.

At a selected contact belonging to \((i,j)\),
the weight \(q_{ij}\) has valuation \(-1\).
Every other weight in the finite global
sum is integral at that prime: its
distinct fixed contact factors and
denominators are units by the prime
selection and avoidance conditions.
It follows that
\[
                    v_p(c)=-1,\qquad
                    s_{ij}=c-q_{ij}\in\mathbb Z_p .
\]
The block \(i\) detector consequently has
exactly its former successful-slot
behavior: \(W_i\equiv1\pmod p\) there,
while at the other slot every factor
is either a unit of residue one or
has a pole. If the other \(W_i\) value
is a unit, its trace is 2 modulo \(p\)
and the shifted norm has positive
valuation. If it has pole order \(r\),
the trace has valuation \(-r\) and the
shifted norm has valuation at least
\(1-r\). In both cases
\[
                  v_p(D_i)>0,\qquad
                  g_{ij}\equiv1\pmod p,\qquad
                  v_p(b_{ij})=1 .
\]
The case \(D_i=0\) is included.

It is unnecessary for \(W_l\), with \(l\ne i\),
to be successful at this prime. For any
foreign pole input \(b_{lr}\), the source
contact factorization gives even valuation
of \(R_{lr}\) away from its own selected
support. Formula \eqref{eq:global-pole}
changes that valuation only by the
valuation of a rational square. It
therefore remains even, even if \(c\),
\(D_l\), or \(g_{lr}\) has a pole.
Foreign contacts cannot create a negative
odd denominator valuation. Combining
this parity statement with the positive
odd valuations at each pole's own
selected contacts gives the same
square-reduced-denominator property
as in the source.

At every prepared finite place and at
real infinity the reflected contact
functions have the common leading
asymptotic
\[
 q_{ij}(\lambda+\delta_i+u_{ij})
             =\kappa\lambda^{-2}+O(\lambda^{-3}),
                         \qquad\kappa>0,
\]
where all shifts have already been fixed.
For the fixed total number \(m=ke\)
of contact functions,
\[
 c=m\kappa\lambda^{-2}+O(\lambda^{-3}),\qquad
 \frac{(c-q_{ij})^2}{q_{ij}}
       =(m-1)^2\kappa\lambda^{-2}+O(\lambda^{-3})
       \longrightarrow0 .
\]
All physical tests are bounded on the
chosen neighborhoods, and each nonconstant
detector term contains a positive power
of \(c\). Thus
\[
          W_i\longrightarrow1,\qquad
          D_i\longrightarrow0,\qquad
          g_{ij}\longrightarrow1 .
\]
These are the required limits also at
two; no division of two vanishing
quantities is being left unjustified.
No division by \(m\) or \(m-1\) is used.
If a prepared prime divides either fixed
integer, a leading coefficient can have
greater valuation, which does not affect
any of the stated limits.
Fixed finite divisibility in the prime
supplier and an arbitrarily large real
slope make all the limits sufficiently
close simultaneously.

The actual native-pole inputs retain the
source local square classes at \(2,3,17\)
and tend at the real place to
\((204/409)^2\in(0,17)\). For sufficiently
large real slope all weights are positive,
so \(c>0\), and the signed factors remain
positive on the prepared neighborhoods.
Apply the retained native realization
theorem to each actual comparison pole
and each scalar pole independently.
Finally fill all point charts and helper
registers using the common nonzero
gauge \(c\). This proves simultaneous
completeness without revisiting any
earlier root or contact choice.

\begin{theorem}[The global contact gauge]
\label{thm:global-gauge}
Under the same retained arithmetic
interfaces and fixed defining parameter
as the common-point construction, replacing
the individual contact sums by
\eqref{eq:global-contact-sum}, with private
\(W_i,D_i\) and pole witnesses, preserves
the rational input projection
\(\mathbb Z^k\). No new prime-selection,
height, or pole-realization hypothesis
is required.
\end{theorem}
\begin{proof}
Soundness follows from
Lemma~\ref{lem:global-gate} and the
unchanged source extraction and height
argument. The ordered construction above
gives a rational witness for each integer
tuple. The literal guard \(c\,c^{-1}=1\)
is retained in the compiled system.
\end{proof}

\section{Removing a private field from every input pair}
\label{sec:global-input-pair}

The shared gauge now permits a quadratic
chart reduction which was unavailable
when the gauges differed between input
blocks. Let
\[
      s=y(U)/x(U),\qquad F=1/x(U)
\]
be the common input-numerator fields.
For block \(i\), the original physical
finite-pair chart is
\[
\begin{split}
 d_i&=F(\eta_i+s+t_i)=\eta_iG_i,\\
 D(F+G_i)&=(s-t_i)\eta_i-t_i^2,\\
 1&=d_i(s-t_i)-DFG_i,
                         \qquad D=11-6\sqrt2 .
\end{split}
\]
Its stored scales in both interfaces are
\[
             G_i=c\,\widetilde G_i,\qquad
             d_i=c\,\widetilde d_i,
\]
while \(s,F,t_i,\eta_i\) have contact
weight zero. The two rational coordinates
of \(\widetilde d_i\) occur only in this
chart.

Introduce two common fields \(x,\widehat F\)
and impose
\begin{equation}\label{eq:global-common-U}
               F=c\widehat F,\qquad
               Fx=1,\qquad x=s^2+DF .
\end{equation}
Replace every input pair by just the
two field equations
\begin{equation}\label{eq:global-short-pair}
 \widehat F(\eta_i+s+t_i)=\eta_i\widetilde G_i,\qquad
 D(F+c\widetilde G_i)=(s-t_i)\eta_i-t_i^2 .
\end{equation}
All monomials have degree at most two.

\begin{lemma}[Exact chart projection]\label{lem:global-short-pair}
Equations \eqref{eq:global-common-U} and
\eqref{eq:global-short-pair} have exactly
the same projection to retained variables
as the original \(k\) charts with common
\(s,F,c\ne0\). They save \(2k-4\)
rational coordinates and \(4k-6\)
rational quadratic residuals.
\end{lemma}
\begin{proof}
For one original chart put
\[
 r_1=d-F(\eta+s+t),\quad
 r_3=D(F+G)-(s-t)\eta+t^2,\quad
 r_4=d(s-t)-DFG-1.
\]
Direct expansion gives
\[
                  Fs^2+DF^2-1
                   =r_4-(s-t)r_1+Fr_3 .
\]
An original solution therefore has
\(F\ne0\), and the unique values
\(x=F^{-1}\), \(\widehat F=F/c\)
satisfy all the new equations.

Conversely define
\(\widetilde d_i=\widehat F(\eta_i+s+t_i)\),
\(G_i=c\widetilde G_i\), and
\(d_i=c\widetilde d_i\).
The first two old chart relations follow
from \eqref{eq:global-short-pair} and
\(F=c\widehat F\). The common curve
relations give \(Fs^2+DF^2=1\).
The displayed identity therefore supplies
the old unit relation. Every original
chart is recovered, including its
tangent and zero-\(d_i\) boundaries.
No division by a private chart field
has been imposed.

We remove one field per input and add
two common fields, a saving of \(2k-4\)
rational coordinates. The \(4k\)
original field equations become
\(2k+3\), giving \(4k-6\) fewer
rational component equations.
For \(k=12\) these savings are
20 coordinates and 42 residuals.
\end{proof}

This algebraic replacement changes no
detector, contact function, or pole
equation. Its exact identity and
source-coordinate occurrence obligations
are also recorded in the independent
checker
\nolinkurl{undecidability/check_point_chart_lemmas.py}.
\section{The common physical functions and their stored coordinates}
\label{sec:global-dictionary}

We now identify every coordinate counted as common. This is essential:
equality of polynomial row names alone does not justify identifying
independent witnesses. The common reference \(H\), the common numerator
\(U_1\), and the common point roles 2, 3, and 4 are supplied by the
ordered construction proved in \cite{common,crt}. The new theorem above
also makes their contact gauge \(c\) common. The input differences
\(Q_{1,i}\), their alternate points \(V_{1,i}\), the scalar normalizers
\(R_i\), the contact arguments, and the native-pole witnesses remain
private.

Write
\[
 f=\frac1{x(H)},\qquad S=\frac{y(H)}{x(H)},\qquad
 T_H=\frac{y(H-P_0)}{x(H-P_0)},\qquad R_H=f(H-P_0)^2.
\]
For the common finite pairs in roles 3 and 4, and for the common input
numerator in role 1, use the fields \(s_j,t_j,F_j,G_j,\eta_j,d_j\)
defined by the finite-pair equations of the preceding section.
Thus \(F_1,s_1\) are common, while \(G_{1,i},t_{1,i},\eta_{1,i}\)
remain private. For role 2, retain the normalized inverse-Vieta fields
\(U,V,W,Z,r,r_2\), where
\[
 Z=2z(Q_2)/z(H),\qquad
 (U,V,W)=Z(u,v,w),\qquad r=Z\rho,\qquad r_2=r^2.
\]
The role-2 ratios and reciprocal squares are
\(A=L_2\), \(F_{U_2}=f(U_2)^2\), and
\(F_{V_2}=f(V_2)^2\). None of these actual physical functions depends
on the private scalar normalizer \(R_i\).

Put \(M=-2S-W\). The common comparison fields are
\[
 E_2^*=A\overline{V+M}+2\overline{T_H}\,\overline Z,
 \quad E_3=S\overline{t_3}-\overline{T_H}s_3,
 \quad E_4=S\overline{t_4}-\overline{T_H}s_4.
\]
For the Lucas circuit write
\[
 \mathcal C=(A-1)(s_4+4S),\qquad
 \mathcal P s_3=2Ss_4,\qquad
 \mathcal Q=(A+1)\mathcal P,
 \qquad
 \mathcal H=\mathcal P^2+\mathcal C^2-\mathcal C\mathcal Q-s_4^2.
\]
The completeness construction chooses the finite neighborhoods where
\(s_3\ne0\), so the same \(\mathcal P\) fills every block. No new
inverse of \(s_3\) is imposed on an arbitrary rational solution.

The four mixed square fields now shared are
\[
 E_3^2-c^3(E_2^*)^2,\quad E_3^2+c^3(E_2^*)^2,
 \quad E_4^2+c\mathcal H^2,\quad E_4^2-c\mathcal H^2.
\]
Their products supply exactly the original enhanced comparison terms.
The three common detector factors, before the unchanged stored
normalizations, are
\begin{align*}
 \Phi_1&=1+c^2F_3^2+c^4\overline{F_{V_2}}^{\,2}
                +c^5E_4^4-c^7\mathcal H^4,\\
 \Phi_2&=1+c^2\overline{G_4}^{\,2}+c^4F_{U_2}^2
                +c^5E_3^4-c^{11}(E_2^*)^4,\\
 \Phi_3&=1+cF_1^2+c^2\overline{G_3}^{\,2}
                -c^3F_4^2+c^4\overline{R_H}^{\,2}.
\end{align*}
The private factor \(\Phi_{0,i}\), containing the input comparison,
is unchanged. In particular no detector coefficient, sign, or
valuation profile has been altered to obtain these savings.

Let \(X_j\) denote the literal physical source coordinates, and
\(Y_j\) their stored coordinates before final renumbering. Their maps
are
\begin{equation}\label{eq:global-stored-map}
                 X_j=c^{\omega_j}R_i^{\rho_j}C^{\kappa_j}Y_j.
\end{equation}
Here \(C\) is fixed and common, and equals 1 in the specified interface.
The source's fixed factor of two in certain imaginary components is
retained. The following complete dictionary lists pairs of rational
coordinates; hence each field contributes two. All listed coordinates
have \(\rho_j=0\), except that the pair \(190:191\) is omitted in
the 163-coordinate interface because its stored \(R_i\)-weight is
\(-1\). Common powers of \(c\) and \(C\) do not obstruct sharing.

\begin{center}\small
\begin{longtable}{@{}p{67mm}p{58mm}r r@{}}
\toprule
Actual common functions & Source pairs & 163 & 179\\
\midrule\endfirsthead
\toprule
Actual common functions & Source pairs & 163 & 179\\
\midrule\endhead
\(f,S,R_H,f^2,T_H\) & \(15:16,17:18,19:20,64:65,150:151\) &10&10\\
\(A,Z,U,V,W,F_{U_2},F_{V_2},r_2\) &
\(\begin{gathered}33:34,102:103,182:187,\\54:55,23:24,194:195\end{gathered}\)&16&16\\
\(r\), where its stored normalizer weight vanishes & \(190:191\)&0&2\\
\(E_2^*,E_3,E_4\)&\(35:36,39:40,43:44\)&6&6\\
\(\mathcal P,\mathcal C,\mathcal Q,\mathcal H\)&\(41:42,46:47,48:49,50:51\)&8&8\\
The four mixed square fields displayed above & \(66:73\)&8&8\\
All six finite-pair fields of role 3 & \(213:224\)&12&12\\
All six finite-pair fields of role 4 & \(225:236\)&12&12\\
\(s_1,F_1\), the common input numerator &\(201:202,205:206\)&4&4\\
\(F_1+cF_4\)&\(610:611\)&2&2\\
\(\Phi_1,\Phi_2,\Phi_3\), with their original fixed normalizations
&\(120:125\)&6&6\\
The two reordered products defined below &\(128:131\)&4&4\\
\(c,c^{-1},c^{-2}\)&\(400,401,402\)&3&3\\
\midrule
Total common rational coordinates&&91&93\\
\bottomrule
\end{longtable}
\end{center}

For clarity, some of these functions are not equal to the raw source
pair. For example
\[
 X_{33}+\sqrt2X_{34}=cA/C,\qquad
 X_{35}+\sqrt2X_{36}=cE_2^*,\qquad
 X_{50}+2\sqrt2X_{51}=c\mathcal H.
\]
The fields \(F_{U_2},F_{V_2},R_H\) likewise retain their fixed
imaginary-component convention. Once \(c\) and \(C\) are common,
the entire scale in \eqref{eq:global-stored-map} is common precisely
for the coordinates counted above. The verifier checks every relevant
\(R_i\)-weight directly against the source map. The private fields
\(R_i x(Q_2)\), the scalar trace data, and the private input alternate
point are excluded from the dictionary.

\section{Reordering the common detector products}

Let \(f_0,f_1,f_2,f_3\) denote the stored field factors, with their
source normalizations already included, and put \(v=c^{-2}\).
Both source interfaces have the following stored multiplication
chain:
\[
                vA_i=f_{0,i}f_1,\qquad
                vB_i=A_i f_2,\qquad
                vW_i=B_i f_3.
\]
The registers \(A_i,B_i\), source pairs \(128:129\) and
\(130:131\), occur only in these equations. Replace them by the
common registers \(A_*,B_*\) defined by
\begin{equation}\label{eq:global-reordered-product}
                vA_*=f_1f_2,\qquad
                vB_*=A_*f_3,\qquad
                vW_i=f_{0,i}B_*.
\end{equation}
All equations remain quadratic over \(K\), and hence give two
rational quadratics each. Because \(c\ne0\), both chains yield
exactly
\[
                       W_i=c^6f_{0,i}f_1f_2f_3.
\]
Conversely their intermediate registers are uniquely recovered by
the displayed products. This is an exact change of auxiliary
coordinates; it does not identify the different \(W_i\).
It saves four rational coordinates per additional input block.
The compiler's occurrence check excludes all other uses of the
redefined registers.

\section{Compiling the global contact sum at quadratic degree}

Every old block stored all but its last contact weight. The last was
the difference between its local gauge and the preceding weights.
Introduce one independent physical last weight \(Q_i\) per block.
Its stored variable is \(Y_{i,*}=Q_i/c^2\), the same contact weight
as the other stored contact weights. The only three local equations
to be changed are the last balanced pole, inverse, and gate equations.
In physical variables they become
\begin{align}
 167281\,c\,b_i(1+Q_i)-166464\,c^2n_i^2&=0,\label{eq:last-pole}\\
 c-Q_i\iota_i&=0,\label{eq:last-inverse}\\
 2cn_i-1-D_i^2-c\iota_i+2c-Q_i&=0.\label{eq:last-gate}
\end{align}
The \(n_i,D_i,\iota_i\) here denote physical gate coordinates,
unrelated to integer point indices. Substituting the old expression
for \(Q_i\) recovers the three old physical residuals identically.
The exact existing Laurent row scales transform these equations into
ordinary quadratics. The independent verifier checks those three
physical identities and all stored monomials directly.

Write \(u=c^{-1}\) and \(v=c^{-2}\), and let \(Y\) be the new
last stored weight. For the 163-coordinate interface the stored rows
are
\begin{align*}
 167281 B(v+Y)-166464 N^2&=0,\\
 uv-IY&=0,\\
 2Nv-E^2-Iu+2uv-v^2-vY&=0.
\end{align*}
For the 179-coordinate interface they are
\begin{align*}
 167281 B(v+Y)-166464 N^2&=0,\\
 uv-IY&=0,\\
 2Nu-E^2-Ic+2u-v-Y&=0.
\end{align*}
The letters \(B,N,I,E\) in these displays are the exact stored
coordinates specified by the table maps. Their different scales
account for the difference between the two third rows.

Since every physical weight is \(c^2Y_{ij}\), the global sum becomes
the linear equation
\begin{equation}\label{eq:linear-weight-sum}
                         \sum_{i,j}Y_{ij}=u.
\end{equation}
Eliminate the last new weight from this equation by affine
substitution. There are twelve newly introduced last weights and one
is eliminated, so the cost is eleven rational coordinates and no
extra residual. The inverse equation \eqref{eq:last-inverse}
continues to force that last weight to be nonzero. No hidden
division by the eliminated affine expression is performed.

\section{Exact counts and the resulting quartics}

The starting coordinate counts are 1601 and 1837. They already shared
54 and 50 rational coordinates, respectively. The common dictionary
increases these to 91 and 93, saving \(11(91-54)=407\) and
\(11(93-50)=473\) further coordinates. The global contact sum costs
eleven coordinates; the common-input chart removes twenty. Thus
\[
       1601-407+11-20=1185,\qquad
       1837-473+11-20=1355.
\]
Independently, from the original 251-coordinate outer circuit the
same counts are
\[
       251+12(m-1)-11s+11-24+4,
       \qquad (m,s)=(163,91)\ \hbox{or}\ (179,93).
\]
The chart replacement removes 42 quadratic residuals. After all
identifications and product reordering, there are respectively 462
and 528 literal duplicate residuals. Consequently
\[
              1599-42-462=1095,\qquad
              1811-42-528=1241.
\]
Each deleted duplicate is listed with the name of its retained
identical equation; these are literal polynomial identities, not
an assertion about generic rank or a heuristic equation count.

\begin{theorem}[Conditional quartic transfer]
Retain the stated arithmetic interfaces and the selected defining
parameters. The complete tables constructed here have 1185 or 1355
rational coordinates and 1095 or 1241 quadratic residuals. Their
outer input projection is exactly that of the frozen emended
positive-integer source, on its retained cofinite parameter set.
Their sums of squares have degree exactly four.

If, in addition, \(U_{\mathrm{em}}\) holds for that exact frozen
integer source, these families yield the corresponding undecidability
reductions. This additional universality assertion remains unresolved.
No unconditional undecidability bound or global minimum follows from
the finite verifications.
\end{theorem}
\begin{proof}
Given a rational solution, expand all shared coordinates into their
separate block copies. Recover the old intermediate product registers
from their unique products, recover every deleted input-chart field
by the exact chart identity, and recover the eliminated weight by
\eqref{eq:linear-weight-sum}. The global-gauge soundness proof applies
to each block independently and makes every distinguished input
integral. This uses the new gate proof, not a claim that the former
local contact sums are restored. The unchanged outer reconstruction
then gives a positive-integer witness for the exact emended source.

Conversely, start from such a positive-integer witness. The common
reference and common numerator preparation applies to its finite
tuple of integral inputs. The global-gauge completeness proof fills
the private contacts, products, and poles. Every function in the
dictionary is genuinely common on those witnesses. The exact chart
replacement and reordered products therefore fill the complete new
system. This proves equality of projections without
\(U_{\mathrm{em}}\); that assertion is invoked only for the
undecidability consequence.

All residuals are ordinary quadratic polynomials with integer
coefficient polynomials in the declared fixed parameters. A sum of
rational squares vanishes exactly when each residual does. At least
one retained residual, \(cu-1\), has nonzero quadratic part. The sum
of squares of the leading quadratic forms cannot vanish identically
over the reals, so the resulting polynomial has degree exactly four.
\end{proof}

\paragraph{Complete reproducible data.}
The producer and independent verifier are
\begin{center}\small
\nolinkurl{undecidability/code/build_global_gauge.py},\\
\nolinkurl{undecidability/verify_global_gauge.py}.
\end{center}
The complete sparse residual tables and certificates are in
\begin{center}\small
\nolinkurl{undecidability/coefficients/jones1980_global_gauge/}.
\end{center}
They contain every variable name, every residual monomial and
coefficient, the source maps, the coordinate removals, every duplicate
row, and the predecessor and table hashes. The verifier independently
reconstructs each retained and replacement residual; checks the exact
physical gate substitutions and the common-field identity; audits the
complete occurrences of deleted and redefined fields; checks degree
and full variable support; and compares the complete output bytes.
The mathematical status fields are also mandatory checks. Run
\begin{center}\small
\texttt{python undecidability/verify\_global\_gauge.py}
\end{center}
from the repository root. The verifier is read-only.

\begin{thebibliography}{9}
\bibitem{outer} ZQ Diophantine Project,
\emph{The emended Jones outer system and its quartic transfer},
\nolinkurl{undecidability/papers/jones1980_quartic_transfer.tex},
23 November 2026; especially the source reconciliation.
\bibitem{common} ZQ Diophantine Project,
\emph{Common physical points in a combined quadratic system},
\nolinkurl{undecidability/papers/common_points_quartic_checkpoint.tex},
23 November 2026.
\bibitem{crt} ZQ Diophantine Project,
\emph{A common numerator point for simultaneous integer certificates},
\nolinkurl{undecidability/papers/common_numerator_transfers.tex},
23 November 2026; the all-sign order assertion and fixed-modulus
compatibility, including avoidance of seven.
\bibitem{gate} ZQ Diophantine Project, based on the supplied
constructions of Gabe Conant with Astra 6.2,
\emph{Audit and extension of the shared-product quartics},
\nolinkurl{papers/quartic/shared_product_audit_20261119_v78.tex},
19 November 2026.
\bibitem{scalar} ZQ Diophantine Project,
\emph{Scalar and point reductions to 163 and 179 coordinates},
\nolinkurl{papers/quartic/scalar_point163_179_20261120_v82.tex},
20 November 2026; independently reconstructed physical equations in
\nolinkurl{computation/variable_reductions_v82/general/check.py}.
\bibitem{pair} ZQ Diophantine Project,
\emph{An exact input-pair chart reduction and its contact-scale obstruction},
\nolinkurl{undecidability/papers/point_chart_research.tex},
23 November 2026.
\end{thebibliography}
\end{document}
